% Source: https://artofproblemsolving.com/community/c3h2391786_rate_problem_in_algebra
% Stable identity: aops-topic-2391786-q01
Let the distance be $d$ miles and let $T$ hours be the travel time that gets him there exactly on time. Three minutes is $\dfrac{1}{20}$ hour, so
\[\dfrac{d}{40}=T+\dfrac{1}{20},\qquad\dfrac{d}{60}=T-\dfrac{1}{20}.\]
Subtracting, $\dfrac{d}{40}-\dfrac{d}{60}=\dfrac{d}{120}=\dfrac{1}{10}$, so $d=12$ miles, and $T=\dfrac{12}{40}-\dfrac{1}{20}=\dfrac14$ hour, i.e. $15$ minutes.

\textbf{(a)} At $50$ miles per hour the trip takes $\dfrac{12}{50}$ hour $=14.4$ minutes, so he arrives $0.6$ minutes ($36$ seconds) early, not on time.

\textbf{(b)} He needs $\dfrac{12}{\dfrac14}=\boxed{48}$ miles per hour.

The reason is that the travel time, not the speed, is what lies halfway: being $3$ minutes late and $3$ minutes early means that the on-time travel time $T$ is the average of the two travel times $\dfrac{d}{40}$ and $\dfrac{d}{60}$. Time is inversely proportional to speed, so the right speed $v$ satisfies $\dfrac dv=\dfrac12\left(\dfrac{d}{40}+\dfrac{d}{60}\right)$, i.e.
\[v=\dfrac{2}{\dfrac{1}{40}+\dfrac{1}{60}}=\dfrac{2\cdot40\cdot60}{40+60}=48,\]
the harmonic mean of $40$ and $60$, which is less than their arithmetic mean $50$.
